\documentclass{article}
\usepackage{amsmath,amssymb}
\usepackage[a4paper,margin=1in]{geometry}

\begin{document}

\section{$\overline{A \circ B} = \bar{A} \bullet \hat{B}$}

\begin{subequations}

\begin{align}
\overline{A \circ B} &= \overline{(A \ominus B) \oplus B} \\
&= \overline{A \ominus B} \ominus \hat{B} \\
&= (\bar{A} \oplus \hat{B}) \ominus \hat{B} \\
&= \bar{A} \bullet \hat{B}
\end{align}

\end{subequations}

\section{$\overline{A \oplus B} = \bar{A} \ominus \hat{B}$} 

By definition:
\begin{align}
A \oplus B &= \{\mathbf{x} | (\hat{B})_{\mathbf{x}} \cap A \neq \phi\} \\
A \ominus B &= \{\mathbf{w} | B_{\mathbf{w}} \subseteq A\}
\end{align}

\begin{subequations}

Taking the complement of the dilation:
\begin{align}
\overline{A \oplus B} &= \overline{\{\mathbf{x} | (\hat{B})_{\mathbf{x}} \cap A \neq \phi\}} \\
&= \{\mathbf{x} | (\hat{B})_{\mathbf{x}} \cap A = \phi\}
\end{align}

If $(\hat{B})_{\mathbf{x}} \cap A = \phi$, then $(\hat{B})_{\mathbf{x}} \subseteq \bar{A}$

Therefore:
\begin{align}
\overline{A \oplus B} &= \{\mathbf{x} | (\hat{B})_{\mathbf{x}} \cap A = \phi\} \\
&= \{\mathbf{x} | (\hat{B})_{\mathbf{x}} \subseteq \bar{A}\} \\
&= \bar{A} \ominus \hat{B}
\end{align}

\end{subequations}

\end{document}